% Source-derived chapter 07: Fully faithfulness
% Original source: tates-thesis-de-rham-setting
\chapter{Fully faithfulness}
\label{tates-thesis-de-rham-setting-chapter-07}
\source{tates-thesis-de-rham-setting}

\begin{remark}[Source section]
\label{tates-thesis-de-rham-setting-chapter-07-source}
\source{tates-thesis-de-rham-setting}
\source{tates-thesis-de-rham-setting}
This chapter reproduces the corresponding section of the retrieved
LaTeX source, including its definitions, statements, examples, and
proofs. It has not yet been translated into Lean.
\notready
\end{remark}

% The source section title was: Fully faithfulness
\label{s:ff}
\source{tates-thesis-de-rham-setting}

\subsection{Overview}

\subsubsection{}

The main result of this section is the following.

\begin{prop}
\source{tates-thesis-de-rham-setting}
\notready\label{p:ff}
\source{tates-thesis-de-rham-setting}

For every $n \geq 0$, the functor $\Delta_n$ is fully faithful.

\end{prop}

\begin{rem}
\source{tates-thesis-de-rham-setting}
\notready\label{r:kash-y}
\source{tates-thesis-de-rham-setting}

By Lemma \ref{l:lambda-ff}, the pushforward
functors $\IndCoh^*(\sY_n) \to \IndCoh^*(\sY)$ are fully faithful
for $n>0$. Therefore, for $n \neq 0$, the above result is equivalent
to fully faithfulness of $\ol{\Delta}_n$. 

\end{rem}

\begin{rem}
\source{tates-thesis-de-rham-setting}
\notready

As $\lambda_{n,*}^{\IndCoh}(\sF_n)$ is compact in 
$\IndCoh^*(\sY)$, Proposition \ref{p:ff} is equivalent to showing
that the map:
%
\[
W_n^{op} \to \ul{\End}_{\IndCoh^*(\sY)}(\lambda_{n,*}^{\IndCoh}\sF_n)
\]

\noindent is an isomorphism.

\end{rem}

Our proof is essentially by induction. We settle the $n = 0$ case
by explicit analysis. We then use local class field theory
(or the Contou-Carr\`ere pairing) to perform the inductive step.

\subsection{Proof for $n = 0$}

We begin with the base case of our induction.

\subsubsection{A preliminary calculation}

We begin with the following explicit calculation.

Recall the scheme $C = \Spec(k[x,y]/xy)$ from \S \ref{ss:refinement}.
We equip it with the $\bG_m$-action of \S \ref{ss:z-refinement},
so $\deg(x) = 1$ and $\deg(y) = 0$ for the corresponding grading.

Let $\sO_{\bA_x^1/\bG_m} \in \QCoh(C/\bG_m)^{\heart}$ denote the structure
sheaf of the $x$-axis, i.e., the object corresponding to the graded
$k[x,y]/xy$-module $k[x] = k[x,y]/y$ (where the generator has degree $0$).

\begin{lem}
\source{tates-thesis-de-rham-setting}
\notready\label{l:c-ext}
\source{tates-thesis-de-rham-setting}

The unit map:
%
\[
k \to \ul{\End}_{\QCoh(C/\bG_m)}(\sO_{\bA_x^1/\bG_m})
\]

\noindent is an isomorphism.

\end{lem}

\begin{proof}

This follows by a standard $\Ext$-calculation that we give here.

Let $A = k[x,y]/xy$. We let $A(n)$ denote $A$ as a graded $A$-module
where the generator is given degree $n$.\footnote{We remark that the
sign conventions here are the same as in \S \ref{sss:z-pic}.}

The complex:
%
\[
\ldots 
\xar{y\cdot -} A(2) \xar{x\cdot -} A(1) \xar{y\cdot -} A(1) \xar{x\cdot -} A \xar{y\cdot -} \underset{\text{coh. deg. } 0} {A} \to 0 \to \ldots 
\]

\noindent gives a graded free resolution of $k[x] = A/y$.

Therefore, we may calculate $\ul{\End}_{A\mod}(k[x])$ as a
graded vector space via the complex:
%
\[
\ldots \to  0 \to \underset{\text{coh. deg. } 0}{k[x]} 
\xar{0} k[x] \xar{x \cdot-} k[x](-1)
\xar{0} k[x](-1) \xar{x \cdot-} k[x](-2) \xar{0} \ldotsplus
\]

\noindent The (graded) degree $0$ component of this complex,
which computes: $$\ul{\End}_{A\mod^{\bG_m,w}}(k[x]) = 
\ul{\End}_{\QCoh(C/\bG_m)}(\sO_{\bA_x^1/\bG_m}),$$ is:
%
\[
\ldots \to 0 \to \underset{\text{coh. deg. } 0}{k} 
\xar{0} k \xar{x \cdot-} k\cdot x
\xar{0} k\cdot x \xar{x \cdot-} k\cdot x^2 \xar{0} \ldotsplus 
\]

\noindent Clearly this complex is acyclic outside of degree cohomological $0$,
and its $H^0$ is 1-dimensional and generated by the unit. 

\end{proof}

\subsubsection{}

We have the following result, explicitly describing the geometry of our situation.

\begin{lem}
\source{tates-thesis-de-rham-setting}
\notready\label{l:z0-emb}
\source{tates-thesis-de-rham-setting}

\begin{enumerate}

\item There is a canonical isomorphism $\sZ^{\leq 0} \simeq \bA^1/\bG_m$.
Explicitly, given $(\sL,\nabla,s) \in \sZ^{\leq 0}$, we
take the line $\sL|_0$ with its section $s|_0$ (for $0 \in \sD$ the base-point).

\item The canonical map $\sZ^{\leq 0} \to \sY$ is an ind-closed embedding.
Its formal completion $\sY_{\sZ^{\leq 0}}^{\wedge}$ 
is isomorphic to $(C/\bG_m)_{\bA^1/\bG_m}^{\wedge}$ compatibly with
the above isomorphism $\sZ^{\leq 0} \simeq \bA^1/\bG_m = \bA_x^1/\bG_m$.

\end{enumerate}

\end{lem}

\begin{proof}

We fix a coordinate $t$ and then consider $\LocSys_{\bG_m}$ 
as mapping to $\bB \bG_m$ via the isomorphism
Proposition \ref{p:locsys}. 
For a prestack $S$ over $\LocSys_{\bG_m}$, let $\widetilde{S}$ denote
the base-change $S \times_{\bB \bG_m} \Spec(k)$. (We remark that this
notation is compatible with that of \S \ref{sss:z-pic}.)

By Proposition \ref{p:zn-cl} and Theorem \ref{t:cl}, 
$\sZ^{\leq 0}$ and $\sY$ are classical. Moreover, the formal completion
$\sY_{\sZ^{\leq 0}}^{\wedge}$ is classical: this follows immediately
from Proposition \ref{p:z-flat-axes} \eqref{i:z-axes-1}.

We can explicitly calculate $\widetilde{\sY_{\sZ^{\leq 0}}^{\wedge}}$
by classicalness and Proposition \ref{p:locsys}: it parametrizes
$y \in \bA_0^{1,\wedge}$ (defining the local system $(\sO,d-y\frac{dt}{t})$)
and an element $f \in (\fL \bA^1)_{\fL^+ \bA^1}^{\wedge}$
such that $df = y f \frac{dt}{t}$. 
If we expand  $f = \sum a_i t^i$, we find $i a_i = y a_i$. As $y$ is
nilpotent, this implies $a_i = 0$ for all $i \neq 0$ and $y a_0 = 0$.
Clearly $\widetilde{\sZ^{\leq 0}}$ corresponds to the locus $y = 0$.
Therefore, writing $x$ for $a_0$ and noting that the relevant $\bG_m$-action
scales $x$, we obtain the claims.

\end{proof}

\subsubsection{}

Now the $n = 0$ case of Proposition \ref{p:ff} asserts that the unit map:
%
\[
k \to \ul{\End}_{\IndCoh^*(\sY)}(\sF_0)
\]

\noindent is an isomorphism. Note that $\sF_0$ is just the pushforward
of the structure sheaf of $\sZ^{\leq 0}$ to $\sY$. The calculation of
the above depends only on the formal completion, so we
obtain the result from Lemmas \ref{l:c-ext} and \ref{l:z0-emb}.

\subsection{A fully faithfulness result}

Before proceeding to our induction, we will need the following observations.

\subsubsection{}

Recall the maps $\zeta$ and $\widetilde{\zeta}$ from \S \ref{ss:y-not}.

\begin{prop}
\source{tates-thesis-de-rham-setting}
\notready\label{p:zeta}
\source{tates-thesis-de-rham-setting}

For any $n > 0$, the map:
%
\[
\ul{\End}_{\IndCoh^*(\LocSys_{\bG_m}^{\leq n})}(
\widetilde{\pi}_{n,*}^{\IndCoh}\sO_{\LocSys_{\bG_m}^{\leq n}}) \to 
\ul{\End}_{\IndCoh^*(\sY^{\leq n})}(
\zeta_*^{\IndCoh}\widetilde{\pi}_{n,*}^{\IndCoh}
\sO_{\LocSys_{\bG_m,\log}^{\leq n}}) 
\]

\noindent is an isomorphism.

\end{prop}

We prove this result in \S \ref{ss:zeta-pf}, after some 
preliminary remarks.

\subsubsection{}

Let $A = \oplus_{i \geq 0} A_i $ be a 
$\bZ^{\geq 0}$-graded (classical, say) algebra.

We let $A\mod^{\bG_m,w} \in \DGCat_{cont}$ denote the 
DG category of graded $A$-modules. For $M \in A\mod^{\bG_m,w}$,
we write $M = \oplus_{i \in \bZ} M_i$ for its
decomposition into weight spaces.

\begin{lem}
\source{tates-thesis-de-rham-setting}
\notready\label{l:gr-hom}
\source{tates-thesis-de-rham-setting}

Suppose $M,N \in A\mod^{\bG_m,w}$. Suppose that 
$M$ is concentrated in positive (graded) degrees,
i.e., $M = \oplus_{i>0} M_i$. Suppose $N$ is concentrated
in degree $0$, i.e., $N = N_0$.

Then:
%
\[
\ul{\Hom}_{A\mod^{\bG_m,w}}(M,N) = 0.
\]

\end{lem}

\begin{proof}

Let $A(i) \in A\mod^{\bG_m,w}$ be as in the proof of Lemma \ref{l:c-ext}, i.e.,
$A$ graded with generator in degree $i$. The modules $A(i)$ for $i > 0$ 
generate the subcategory of $A\mod^{\bG_m,w}$ consisting of modules concentrated
in positive degrees. So we reduce to the case $M = A(i)$ for $i>0$, for which the claim is obvious.
	
\end{proof}

\subsubsection{}

Suppose now that we are given a $\bZ^{\geq 0}$-graded 
algebra $B = \oplus_{i \geq 0} B_i$ 
and a graded map $\iota^*:A \to B$ inducing
an isomorphism $A_0 \isom B_0$. We let $\iota_*:B\mod^{\bG_m,w} \to A\mod^{\bG_m,w}$
and $\iota^*:A\mod^{\bG_m,w} \to B\mod^{\bG_m,w}$ denote the induced
adjoint functors.

\begin{prop}
\source{tates-thesis-de-rham-setting}
\notready\label{p:gr-hom}
\source{tates-thesis-de-rham-setting}

In the above setting, suppose $M,N \in B\mod^{\bG_m,w}$ with 
$M$ concentrated in non-negative graded\footnote{As opposed to cohomological.} 
degrees and $N$ concentrated in degree $0$. 

Then the natural map:
%
\[
\ul{\Hom}_{B\mod^{\bG_m,w}}(M,N) \to 
\ul{\Hom}_{A\mod^{\bG_m,w}}(\iota_*M,\iota_* N) 
\]

\noindent is an isomorphism.

\end{prop}

\begin{proof}

\step\label{st:gr-1} 
\source{tates-thesis-de-rham-setting}

Let $\widetilde{M} \in A\mod^{\bG_m,w}$ be concentrated in non-negative graded degrees.
We claim that $\Ker(\widetilde{M} \to \iota_*\iota^* \widetilde{M})$ is concentrated
in (strictly) positive graded degrees.

First, in the case $\widetilde{M} = A$, this map is $\Ker(A \to \iota_*(B)$, in which case
the claim follows as we assumed $A \to B$ an isomorphism in graded degree $0$.

In general, we have:
%
\[
\Ker(\widetilde{M} \to \iota_*\iota^* \widetilde{M}) = \Ker(A \to \iota_*(B))  \underset{A}{\otimes}
\widetilde{M} 
\]

\noindent which is the tensor product of a module in graded degrees $\geq 0$ and one in graded degrees $>0$, so 
is in graded degrees $>0$ (as $A$ is non-negatively graded), as claimed.

\step\label{st:gr-2}
\source{tates-thesis-de-rham-setting}

Next, we claim that $\Ker(\iota^*\iota_* M \to M)$ is concentrated in strictly positive graded degrees.

Clearly $\iota^*\iota_* M = B \otimes_A \iota_* M$ is in degrees $\geq 0$, so it suffices to see that
$\iota^*\iota_* M \to M$ is an isomorphism in degree $0$. We can check this after applying $\iota_*$.
Then we have a retraction:
%
\[
\iota_* M \to \iota_* \iota^*\iota_* M \to \iota_* M,
\]

\noindent so it suffices to see that the first map is an isomorphism in degree $0$.
Taking $\widetilde{M} = \iota_*M$, this follows from Step \ref{st:gr-1}.

\step Applying Lemma \ref{l:gr-hom} to $\Ker(\iota^*\iota_*M \to M)$ and $N$, we see that the map:
%
\[
\ul{\Hom}_{B\mod^{\bG_m,w}}(M,N) \to \ul{\Hom}_{B\mod^{\bG_m,w}}(\iota^*\iota_*M,N) 
\]

\noindent is an isomorphism. The right hand side identifies with $\ul{\Hom}_{A\mod^{\bG_m,w}}(\iota_*M,\iota_*N)$
by adjunction, and the induced map is given by $\iota_*$ functoriality, so we obtain the claim.
 	
\end{proof}

\subsubsection{}

We will apply the above in the following setting.

Let $A_n$ be the graded ring of \S \ref{ss:coords}.
We write $A_n = \oplus_{i \geq 0} A_{n,i}$ for 
its decomposition into graded pieces. We remind
from \emph{loc. cit}. that $\Spec(A_n)/\bG_m = \sZ^{\leq n}$.

By construction, we have:
%
\[
\LocSys_{\bG_m,\log}^{\leq n} \underset{\bB \bG_m}{\times} \Spec(k) = \Spec(A_{n,0})
\]

\noindent such that the graded algebra maps $A_{n,0} \to A_n \to A_{n,0}$
correspond to the maps

\begin{equation}\label{eq:zeta-proj}
\source{tates-thesis-de-rham-setting}
\LocSys_{\bG_m,\log}^{\leq n} \xar{\widetilde{\zeta}} \sZ^{\leq n} \to 
\LocSys_{\bG_m,\log}^{\leq n}
\end{equation}

\noindent (the latter map being the projection).

\subsubsection{}\label{sss:gr-hom}
\source{tates-thesis-de-rham-setting}

Recall the maps $\iota$, $\iota^r$ from \S \ref{ss:y-not}.

The map $\iota$ arises\footnote{Non-canonically: i.e., we need a choice
of coordinate on the disc to turn $\iota$ into a map of 
stacks over $\bB \bG_m$.} from a map $\iota^*:A_n \to A_n$ of graded rings.
This map is an isomorphism
on $0$th graded components, as is evident from 
\eqref{eq:zeta-proj}. 

\begin{cor}
\source{tates-thesis-de-rham-setting}
\notready\label{c:gr-hom}
\source{tates-thesis-de-rham-setting}

Suppose $\sG \in \QCoh(\LocSys_{\bG_m,\log}^{\leq n})$. Then the morphism:
%
\[
\ul{\Hom}_{\QCoh(\sZ^{\leq n})}(\sO_{\sZ^{\leq n}},\widetilde{\zeta}_* \sG) \to 
\ul{\Hom}_{\QCoh(\sZ^{\leq n})}(\iota_*\sO_{\sZ^{\leq n}},\iota_*\widetilde{\zeta}_* \sG) 
\]

\noindent is an isomorphism. More generally, 
for any $r,s \geq 0$, the morphism:
%
\[
\ul{\Hom}_{\QCoh(\sZ^{\leq n})}(\iota_*^r(\sO_{\sZ^{\leq n}}),\iota_*^s\widetilde{\zeta}_* \sG) \to 
\ul{\Hom}_{\QCoh(\sZ^{\leq n})}(\iota_*^{r+1}(\sO_{\sZ^{\leq n}}),\iota_*^{s+1}\widetilde{\zeta}_* \sG) 
\]

\noindent is an isomorphism.
	
\end{cor}

\begin{proof}

Translating to graded modules and using \eqref{eq:zeta-proj},
the hypotheses of Proposition \ref{p:gr-hom} are clearly satisfied, giving the result.
	
\end{proof}

\subsubsection{Conclusion}\label{ss:zeta-pf}
\source{tates-thesis-de-rham-setting}

We now prove Proposition \ref{p:zeta}.

Observe that $\LocSys_{\bG_m}^{\leq n} = \lim_n \sZ^{\leq n}$, where the limit
is formed under the (affine) morphism $\iota$. Therefore, we have:
%
\[
\widetilde{\zeta}_*\sO_{\LocSys_{\bG_m,\log}^{\leq n}} = 
\underset{r}{\colim} \, \iota_*^r \sO_{\sZ^{\leq n}} \in \QCoh(\sZ^{\leq n}).
\]

\noindent As all of these terms are in the heart of the $t$-structure, 
we obtain a similar identity in $\IndCoh(\sZ^{\leq n})$, using $\IndCoh$-pushforwards 
in the notation instead.

Define:
%
\[
\sG \coloneqq \Oblv\Av_!^{\Gr_{\bG_m}^{\leq n},w}
\sO_{\LocSys_{\bG_m,\log}^{\leq n}} \in 
\QCoh(\LocSys_{\bG_m,\log}^{\leq n})^{\heart}.
\]

We let $i:\sZ^{\leq n} \to \sY_{\log}^{\leq n} = \colim_{n,\iota} \sZ^{\leq n}$ 
denote the structural map. Using adjunction, we rewrite:
%
\begin{equation}\label{eq:zeta-rhs}
\source{tates-thesis-de-rham-setting}
\ul{\End}_{\IndCoh^*(\sY^{\leq n})}
(\zeta_*^{\IndCoh}\widetilde{\pi}_{n,*}^{\IndCoh}
\sO_{\LocSys_{\bG_m,\log}^{\leq n}}),
\end{equation}

\noindent which is the right hand side of Proposition \ref{p:zeta}, as:
%
\[
\begin{gathered}
\ul{\Hom}_{\IndCoh^*(\sY_{\log}^{\leq n})}
(i_*^{\IndCoh}\widetilde{\zeta}_*^{\,\IndCoh}\sO_{\LocSys_{\bG_m,\log}^{\leq n}},
i_*^{\IndCoh}\widetilde{\zeta}_*^{\,\IndCoh}\Oblv\Av_!^{\bG_m^{\leq n},w}\sO_{\LocSys_{\bG_m,\log}^{\leq n}}) = \\
\ul{\Hom}_{\IndCoh^*(\sY_{\log}^{\leq n})}
(i_*^{\IndCoh}\widetilde{\zeta}_*^{\,\IndCoh}\sO_{\LocSys_{\bG_m,\log}^{\leq n}},
i_*^{\IndCoh}\widetilde{\zeta}_*^{\,\IndCoh}\sG). 
\end{gathered}
\]

\noindent By the above, we can further express this term as:
%
\[
\begin{gathered}
\ul{\Hom}_{\IndCoh^*(\sY_{\log}^{\leq n})}
(i_*^{\IndCoh}\underset{r}{\colim} \, \iota_*^{r,\IndCoh} \sO_{\sZ^{\leq n}},
i_*^{\IndCoh}\widetilde{\zeta}_*^{\,\IndCoh}\sG) = \\
\underset{r}{\lim} \, \ul{\Hom}_{\IndCoh^*(\sY_{\log}^{\leq n})}
(i_*^{\IndCoh} \iota_*^{r,\IndCoh} \sO_{\sZ^{\leq n}},
i_*^{\IndCoh}\widetilde{\zeta}_*^{\,\IndCoh}\sG).
\end{gathered}
\]

\noindent As $\iota_*^r \sO_{\sZ^{\leq n}} \in \Coh(\sZ^{\leq n})$ for all $r$,
by \cite{cpsii} Corollary 6.5.3, we may calculate the above as:
%
\[
\underset{r}{\lim} \, 
\underset{s}{\colim} \, 
\ul{\Hom}_{\IndCoh^*(\sZ^{\leq n})}
(\iota_*^{r+s,\IndCoh} \sO_{\sZ^{\leq n}},
\iota_*^s\widetilde{\zeta}_*^{\,\IndCoh}\sG).
\]

Now for fixed $r$, Corollary \ref{c:gr-hom} implies that all of the
maps in the above colimit are isomorphisms. Therefore, the above
may be calculated as:
%
\[
\underset{r}{\lim} \, 
\ul{\Hom}_{\IndCoh^*(\sZ^{\leq n})}
(\iota_*^{r,\IndCoh} \sO_{\sZ^{\leq n}},
\widetilde{\zeta}_*^{\,\IndCoh}\sG).
\]

\noindent Moreover, the analysis from \emph{loc. cit}. shows that
$\sO_{\sZ} \to \iota_*^r\sO_{\sZ}$ corresponds to a map of
non-negatively graded modules that is an isomorphism in degree $0$, 
so Lemma \ref{l:gr-hom} shows that the maps is the above
limit are isomorphisms as well. Therefore, we may calculate this
limit as:
%
\begin{equation}\label{eq:no-r-s}
\source{tates-thesis-de-rham-setting}
\ul{\Hom}_{\IndCoh^*(\sZ^{\leq n})}
(\sO_{\sZ^{\leq n}},
\widetilde{\zeta}_*^{\,\IndCoh}\sG).
\end{equation}

\noindent This term is:
%
\[
\begin{gathered}
\Gamma^{\IndCoh}(\sZ^{\leq n},\widetilde{\zeta}_*^{\,\IndCoh}\sG) = 
\Gamma^{\IndCoh}(\LocSys_{\bG_m,\log}^{\leq n},\sG) = \\
\ul{\Hom}_{\IndCoh^*(\LocSys_{\bG_m,\log}^{\leq n})}
(\sO_{\LocSys_{\bG_m,\log}^{\leq n}},\Oblv\Av_!^{\Gr_{\bG_m}^{\leq n},w} 
\sO_{\LocSys_{\bG_m,\log}^{\leq n}}) = \\ 
\ul{\End}_{\IndCoh^*(\LocSys_{\bG_m}^{\leq n})}
(\Av_!^{\Gr_{\bG_m}^{\leq n},w} 
\sO_{\LocSys_{\bG_m,\log}^{\leq n}}).
\end{gathered}
\]

\noindent This last expression is another way of writing the left hand side
of Proposition \ref{p:zeta}. It is straightforward to see that the isomorphism
we have just constructed is inverse to the map in \emph{loc. cit}., yielding the result.


\subsection{Structure of the argument}

We are now positioned to prove Proposition \ref{p:ff}.
We begin by outlining the argument.

\subsubsection{}

The starting point is the following elementary observation.

\begin{lem}
\source{tates-thesis-de-rham-setting}
\notready\label{l:weyl-bimod}
\source{tates-thesis-de-rham-setting}

Suppose $f:M \to N \in W_2\mod$ is a morphism (in the derived category) of 
modules over the Weyl algebra $W_2$ in two variables,
with generators denoted $\alpha$, $\beta$, $\partial_{\alpha}$ and $\partial_{\beta}$. 
Then $f$ is an isomorphism if and only if the morphisms:
%
\begin{equation}\label{eq:weyl-bimod}
\source{tates-thesis-de-rham-setting}
\begin{gathered}
\underset{r,\alpha}{\colim} \, M \coloneqq \colim(M \xar{\alpha} M \xar{\alpha} \ldots) \to 
\underset{r,\alpha}{\colim} \, N \\
\underset{r,\beta}{\colim} \, M \to \underset{r,\beta}{\colim} \, N \\
M^{\alpha}/\beta \to N^{\alpha}/\beta	
\end{gathered}
\end{equation}

\noindent are isomorphisms in $\Vect$. 
In the last line, we define e.g. $M^{\alpha}$ as the
(homotopy) kernel of $\alpha:M \to M$, and the quotient by $\beta$ means
the homotopy cokernel of the induced map $\beta:M^{\alpha} \to M^{\alpha}$
(which exists because $\alpha$ and $\beta$ commute).
	
\end{lem}

\begin{proof}

In terms of $D$-modules on $\bA^2$, our assumptions are that $f$ induces an isomorphism 
on restriction to $(\bA^1\setminus 0) \times \bA^1$, $\bA^1 \times (\bA^1 \setminus 0)$,
and on $!$-restriction to $0$ (up to a cohomological shift). This implies the claim by 
$D$-module excision.

\end{proof}

\subsubsection{}\label{sss:ff-induction-strat}
\source{tates-thesis-de-rham-setting}

We choose coordinates on the disc $\sD$. We obtain an isomorphism 
$\fL_n^+ \bA^1 \simeq \bA^n$; we let $\alpha_0,\ldots,\alpha_{n-1}$ denote the coordinates
(so e.g. $\alpha_0 = \ev$ as maps to $\bA^1$).
We obtain an evident decomposition $W_n \simeq W_1 \otimes W_{n-1}$, where the first
tensor factor uses the coordinate $\alpha_0$.

Let $f$ denote the morphism $W_n^{op} \to \ul{\End}_{\IndCoh^*(\sY^{\leq n})}(\sF_n) \in \Alg$.
As we have a morphism $W_1^{op} \to W_n^{op}$ of algebras as above, we may regard $f$
as a morphism of $W_1$-bimodules (via left and right multiplication).

Our argument applies Lemma \ref{l:weyl-bimod} using this bimodule structure.
Generally speaking, the first two conditions from the lemma 
amount to Proposition \ref{p:zeta} and its relatives,
and the last one is given by induction.

\subsection{Verifying the axioms} % better heading?

\subsubsection{}

In the setting of \S \ref{sss:ff-induction-strat}, 
it remains to verify the conditions of Lemma \ref{l:weyl-bimod}.
We do so below. 

\subsubsection{Colimits}

We begin with checking the first two maps from \eqref{eq:weyl-bimod} are isomorphisms.

Clearly:
%
\[
\underset{r,\alpha_0\cdot-}{\colim} \, W_n = \underset{r,-\cdot \alpha_0}{\colim} \, W_n = W_n[\alpha_0^{-1}]
= \Gamma(\fL_n^+ \bG_m,D_{\fL_n^+ \bG_m}).
\]

\noindent I.e., in our example, the left hand sides of the first two equations in \eqref{eq:weyl-bimod}
each identify with global differential operators on $\fL_n^+ \bG_m$.

Below, we verify the same for the right hand sides in our examples. Then we check that our map ($f$, in the notation of Lemma \ref{l:weyl-bimod}), is compatible
with these identifications.

\subsubsection{}\label{sss:alpha-0-left}
\source{tates-thesis-de-rham-setting}

The left action of $\alpha_0$ on $\ul{\End}_{\IndCoh^*(\sY^{\leq n})}(\sF_n)$ 
is, by construction, obtained by right composition with the map:
%
\[
\sF_n = 
\Av_!^{\Gr_{\bG_m}^{\leq n},w}(i_{n,*}^{\IndCoh}(\sO_{\sZ^{\leq n}})) \to 
\Av_!^{\Gr_{\bG_m}^{\leq n},w}(\iota_*^{\IndCoh}i_{n,*}^{\IndCoh}(\sO_{\sZ^{\leq n}})) = \sF_n
\]

\noindent obtained by applying $\Av_!^{\Gr_{\bG_m}^{\leq n},w}$ to 
the canonical map:
%
\[
\sO_{\sZ^{\leq n}} \to \iota_*^{\IndCoh}\sO_{\sZ^{\leq n}}.
\]

Therefore, by compactness of $\sF_n$, we see that:
%
\[
\underset{r,\alpha_0\cdot-}{\colim} \, \ul{\End}_{\IndCoh^*(\sY^{\leq n})}(\sF_n)
\isom 
\ul{\Hom}_{\IndCoh^*(\sY^{\leq n})}(\sF_n, 
\Av_!^{\Gr_{\bG_m}^{\leq n},w} \underset{r}{\colim} \, 
(\iota_*^{r,\IndCoh})i_{n,*}^{\IndCoh}(\sO_{\sZ^{\leq n}}).
\]

\noindent The right hand side clearly identifies with:
%
\[
\ul{\Hom}_{\IndCoh^*(\sY^{\leq n})}(\sF_n, 
\zeta_*^{\IndCoh}
\Oblv\Av_!^{\Gr_{\bG_m}^{\leq n},w} (\sO_{\LocSys_{\bG_m,\log}^{\leq n}})).
\]

\noindent Unwinding the definition of $\sF_n$, we can further identify it with:
%
\[
\ul{\Hom}_{\IndCoh^*(\sY_{\log}^{\leq n})}\big(i_{n,*}^{\IndCoh}\sO_{\sZ^{\leq n}}, 
i_{n,*}^{\IndCoh}\widetilde{\zeta}_*^{\,\IndCoh}
\Oblv\Av_!^{\Gr_{\bG_m}^{\leq n},w} (\sO_{\LocSys_{\bG_m,\log}^{\leq n}})\big).
\]

\noindent The comparison of \eqref{eq:zeta-rhs} and \eqref{eq:no-r-s} 
from \S \ref{ss:zeta-pf} yields that the map:
%
\[
\begin{gathered}
\Gamma\big(\LocSys_{\bG_m,\log}^{\leq n},\Oblv\Av_!^{\Gr_{\bG_m}^{\leq n},w} (\sO_{\LocSys_{\bG_m,\log}^{\leq n}})\big) = \\
\Gamma^{\IndCoh}\big(\sZ^{\leq n},\widetilde{\zeta}_*^{\,\IndCoh}
\Oblv\Av_!^{\Gr_{\bG_m}^{\leq n},w})(\sO_{\LocSys_{\bG_m,\log}^{\leq n}})\big) = \\
\ul{\Hom}_{\IndCoh^*(\sZ^{\leq n})}\big(\sO_{\sZ^{\leq n}}, 
\widetilde{\zeta}_*^{\,\IndCoh}
\Oblv\Av_!^{\Gr_{\bG_m}^{\leq n},w} (\sO_{\LocSys_{\bG_m,\log}^{\leq n}})\big) \to \\
\ul{\Hom}_{\IndCoh^*(\sY_{\log}^{\leq n})}\big(i_{n,*}^{\IndCoh}\sO_{\sZ^{\leq n}}, 
i_{n,*}^{\IndCoh}\widetilde{\zeta}_*^{\,\IndCoh}
\Oblv\Av_!^{\Gr_{\bG_m}^{\leq n},w} (\sO_{\LocSys_{\bG_m,\log}^{\leq n}})\big)
\end{gathered}
\]

\noindent is an isomorphism. 

By Theorem \ref{t:trun-lcft}, the left hand side identifies canonically with:
%
\[
\Gamma(\fL_n^+ \bG_m,D_{\fL_n^+ \bG_m}).
\]

Putting this together, we obtain an isomorphism:
%
\[
\underset{r,\alpha_0\cdot-}{\colim} \, \ul{\End}_{\IndCoh^*(\sY^{\leq n})}(\sF_n)
\simeq 
\Gamma(\fL_n^+ \bG_m,D_{\fL_n^+ \bG_m}).
\] 

\subsubsection{}\label{sss:alpha-alg}
\source{tates-thesis-de-rham-setting}

In \S \ref{sss:alpha-0-left}, we constructed an isomorphism:
%
\[
\underset{r,\alpha_0\cdot-}{\colim} \, \ul{\End}_{\IndCoh^*(\sY^{\leq n})}(\sF_n)
\simeq 
\Gamma(\fL_n^+ \bG_m,D_{\fL_n^+ \bG_m}).
\]

\noindent We claim that the resulting map:

\[
\ul{\End}_{\IndCoh^*(\sY^{\leq n})}(\sF_n) \to 
\underset{r,\alpha_0\cdot-}{\colim} \, \ul{\End}_{\IndCoh^*(\sY^{\leq n})}(\sF_n)
\simeq \Gamma(\fL_n^+ \bG_m,D_{\fL_n^+ \bG_m})
\]

\noindent is in fact (canonically) a map of (DG) algebras, at least
once the target is given the opposite multiplication. 

For convenience, define:
%
\[
\o{\sF}_n \coloneqq 
\zeta_*^{\IndCoh}\Av_!^{\Gr_{\bG_m}^{\leq n},w}\sO_{\LocSys_{\bG_m,\log}^{\leq n}}.
\]

\noindent Note that there is a canonical morphism:
%
\[
\sF_n \to \o{\sF}_n.
\]

\noindent We in effect showed that the natural maps:
%
\[
\ul{\Hom}_{\IndCoh^*(\sY^{\leq n})}(\sF_n,\o{\sF}_n) \to 
\ul{\End}_{\IndCoh^*(\sY^{\leq n})}(\o{\sF}_n)  \leftarrow 
\ul{\End}_{\IndCoh^*(\LocSys_{\bG_m}^{\leq n})}
(\Av_!^{\Gr_{\bG_m}^{\leq n},w}\sO_{\LocSys_{\bG_m,\log}^{\leq n}})
\]

\noindent are isomorphisms. The map:

\[
\ul{\End}_{\IndCoh^*(\sY^{\leq n})}(\sF_n) \to 
\underset{r,\alpha_0\cdot-}{\colim} \, \ul{\End}_{\IndCoh^*(\sY^{\leq n})}(\sF_n)
\isom 
\ul{\Hom}_{\IndCoh^*(\sY^{\leq n})}(\sF_n,\o{\sF}_n)
\]

\noindent is clearly a map of right 
$\ul{\End}_{\IndCoh^*(\sY^{\leq n})}(\sF_n)$-modules.
But $\ul{\Hom}_{\IndCoh^*(\sY^{\leq n})}(\sF_n,\o{\sF}_n)$ carries
a commuting left 
$\ul{\End}_{\IndCoh^*(\sY^{\leq n})}(\o{\sF}_n)$-module structure,
for which it is a free module of rank $1$ by the above. 
As the class field theory isomorphism:
%
\begin{equation}\label{eq:lcft-diff}
\source{tates-thesis-de-rham-setting}
\ul{\End}_{\IndCoh^*(\LocSys_{\bG_m}^{\leq n})}
(\Av_!^{\Gr_{\bG_m}^{\leq n},w}\sO_{\LocSys_{\bG_m,\log}^{\leq n}})
\simeq \Gamma(\fL_n^+ \bG_m,D_{\fL_n^+ \bG_m})
\end{equation}

\noindent is an isomorphism of algebras, we obtain the claim.

\subsubsection{}\label{sss:lcft-compat-left}
\source{tates-thesis-de-rham-setting}

The above calculations verify that there are canonical isomorphisms:
%
\[
\underset{r,\alpha_0\cdot-}{\colim} \, W_n \simeq 
\Gamma(\fL_n^+ \bG_m,D_{\fL_n^+ \bG_m}) \simeq 
\underset{r,\alpha_0\cdot-}{\colim} \, \ul{\End}_{\IndCoh^*(\sY^{\leq n})}(\sF_n).
\]

We claim that this map is induced by our comparison map $f$ from 
\S \ref{sss:ff-induction-strat}.

By \S \ref{sss:alpha-alg}, it suffices to check that the composition:
%
\[
W_n^{op} \to \ul{\End}_{\IndCoh^*(\sY^{\leq n})}(\sF_n)
\to
\underset{r,\alpha_0\cdot-}{\colim} \, \ul{\End}_{\IndCoh^*(\sY^{\leq n})}(\sF_n)
\simeq \Gamma(\fL_n^+ \bG_m,D_{\fL_n^+ \bG_m})^{op}
\]

\noindent is the restriction map. We first check this on some particular elements.

For the linear functions $t^{-n}k[[t]]dt/k[[t]]dt \subset W_n^{op}$,
this follows by construction. Indeed, the geometric incarnation of the
action of these elements is the observation that 
$\sZ^{\leq n} \subset \sY_{\log}^{\leq n}$ is closed under
the action of $\Gr_{\bG_m}^{pos,\leq n}$. The same is true
for $\LocSys_{\bG_m,\log}^{\leq n} \subset \sY_{\log}^{\leq n}$ \textemdash{}
in fact, it is closed under $\Gr_{\bG_m}^{\leq n}$. Tracing the
constructions, these observations provide the claim.

Next, we claim that the vector fields on $\fL_n^+\bA^1$ defined by the
infinitesimal $\fL_n^+\bG_m$-action: 
%
\[
k[[t]]/t^n \simeq \Lie(\fL_n^+\bG_m) \subset W_n^{op}
\]

\noindent have the correct image. Indeed, this follows by 
construction of the class field theory isomorphism and Lemma \ref{l:hc-calc}.

By the above, we see that the algebra map:

\[
W_n^{op} \to \Gamma(\fL_n^+ \bG_m,D_{\fL_n^+ \bG_m})^{op}
\]

\noindent factors through the (non-commutative) localization of
$W_n^{op}$ at $\alpha_0$, i.e., through a map:

\[
W_n^{op} \to \Gamma(\fL_n^+ \bG_m,D_{\fL_n^+ \bG_m})^{op} \xar{\gamma}
\Gamma(\fL_n^+ \bG_m,D_{\fL_n^+ \bG_m})^{op}
\]

\noindent for some algebra map $\gamma$. As the subspaces
$t^{-n}k[[t]]dt/k[[t]]dt$ and $k[[t]]/t^n$ 
generate $\Gamma(\fL_n^+ \bG_m,D_{\fL_n^+ \bG_m})^{op}$ as an algebra, 
and as we have seen $\gamma$ is the identity on these elements, $\gamma$
itself must be the identity map, as desired.

\subsubsection{}

Next, we calculate:
%
\[
\underset{r,-\cdot \alpha_0}{\colim} \, \ul{\End}_{\IndCoh^*(\sY^{\leq n})}(\sF_n).
\]

For $r \in \bZ$, let $\presup{r}{\sZ^{\leq n}} \subset \sY_{\log}^{\leq n}$
denote the closed where the meromorphic section $s$ has a pole of order $\leq r$;
therefore, $\presup{0}{\sZ^{\leq n}} = \sZ^{\leq n}$,
$\iota(\sZ^{\leq n}) = \presup{-1}{\sZ^{\leq n}}$, and $\alpha_0$ can be
considered as obtained by restriction 
$\presup{r}{\sZ^{\leq n}} \to \presup{r-1}{\sZ^{\leq n}}$ for any $r$.

We let $i_{n,r}:\presup{r}{\sZ^{\leq n}} \to \sY_{\log}^{\leq n}$ denote
the embedding.

For any $\sG \in \IndCoh^*(\sY^{\leq n})$, it follows by definition
that:
%
\[
\underset{r,-\cdot \alpha_0}{\colim} \, \ul{\Hom}_{\IndCoh^*(\sY^{\leq n})}(\sF_n,\sG) =
\underset{r,-\cdot \alpha_0}{\colim} \, \ul{\Hom}_{\IndCoh^*(\sY^{\leq n})}
(i_{n,r,*}^{\IndCoh}(\sO_{\presup{r}\sZ^{\leq n}}),\Oblv\sG).
\]

\noindent We have a canonical isomorphism:
%
\[
\underset{r,-\cdot \alpha_0}{\colim} \, \ul{\Hom}_{\IndCoh^*(\sY^{\leq n})}
(i_{n,r,*}^{\IndCoh}(\sO_{\presup{r}\sZ^{\leq n}}),\Oblv\sG) \isom 
\Gamma^{\IndCoh}(\sY_{\log}^{\leq n},\Oblv \sG).
\]

Now taking $\sG = \sF_n$, we see that:
%
\[
\underset{r,-\cdot \alpha_0}{\colim} \, \ul{\End}_{\IndCoh^*(\sY^{\leq n})}(\sF_n)
\simeq \Gamma^{\IndCoh}(\sY_{\log}^{\leq n},\Oblv \Av_!^{\Gr_{\bG_m}^{\leq n},w}
i_{n,*}^{\IndCoh}\sO_{\sZ^{\leq n}}).
\]

By the analysis of \S \ref{sss:gr-hom},\footnote{Specifically, the 
observations that the ring $A_n$ from 
\emph{loc. cit}. is positively graded with degree $0$ component corresponding to
$\LocSys_{\bG_m,\log}^{\leq n} \subset \sZ^{\leq n} = \Spec(A_n)/\bG_m$.}, 
the map:
%
\[
\Gamma^{\IndCoh}(\sY_{\log}^{\leq n},\Oblv \Av_!^{\Gr_{\bG_m}^{\leq n},w}
i_{n,*}^{\IndCoh}\sO_{\sZ^{\leq n}}) \to 
\Gamma^{\IndCoh}(\sY_{\log}^{\leq n},\Oblv \Av_!^{\Gr_{\bG_m}^{\leq n},w}
i_{n,*}^{\IndCoh}\widetilde{\zeta}_*^{\,\IndCoh}\sO_{\LocSys_{\bG_m,\log}^{\leq n}})
\]

\noindent is an isomorphism. By functoriality, the target identifies with:
%
\[
\Gamma(\LocSys_{\bG_m,\log}^{\leq n},
\Oblv \Av_!^{\Gr_{\bG_m}^{\leq n},w}
\sO_{\LocSys_{\bG_m,\log}^{\leq n}}).
\] 

Now passing to the colimit over $\alpha_0$ acting on the left and
right simultaneously, the above analysis actually shows that the maps: 
%
\[
\underset{r,-\cdot \alpha_0}{\colim} \, \ul{\End}_{\IndCoh^*(\sY^{\leq n})}(\sF_n)
\to \underset{(r_1,r_2),-\cdot \alpha_0,-\cdot \alpha_0}{\colim} \, \ul{\End}_{\IndCoh^*(\sY^{\leq n})}(\sF_n)
\leftarrow 
\underset{r,\alpha_0 \cdot -}{\colim} \, \ul{\End}_{\IndCoh^*(\sY^{\leq n})}(\sF_n.
\]

\noindent are isomorphisms. It then follows from \S \ref{sss:lcft-compat-left} that
the comparison map:

\[
\underset{r,-\cdot \alpha_0}{\colim} \, W_n \to 
\underset{r,-\cdot \alpha_0}{\colim} \, \ul{\End}_{\IndCoh^*(\sY^{\leq n})}(\sF_n)
\]

\noindent is an isomorphism. 

\subsubsection{Hamiltonian reduction}

It remains to check that the third map in \eqref{eq:weyl-bimod} 
is an isomorphism. We will see that this amounts to our
inductive hypothesis.

\subsubsection{}\label{sss:qh-1}
\source{tates-thesis-de-rham-setting}

We begin with some notation.

Suppose $A$ is a (DG) algebra, and that we are given a map:
%
\[
k[t] \to A.
\] 

\noindent We let $\alpha \in \Omega^{\infty} A$ denote the image of $t$.

In this case, we let:
%
\[
\on{QH}_{\alpha}(A) \coloneqq A^{\alpha\cdot-}/-\cdot \alpha 
\]

\noindent with notation on the right hand side understood as
in Lemma \ref{l:weyl-bimod}, i.e., we take the homotopy quotient of 
right multiplication by $\alpha$ and then the homotopy kernel of the
left action of $\alpha$ (and of course, the order of these two operations is
not important).

As is standard, we refer to $\on{QH}_{\alpha}(A)$ 
as the \emph{quantum Hamiltonian reduction} 
of $A$ by $\alpha$.

\subsubsection{}

First, we note that $\on{QH}_{\alpha_0}(W_1^{op}) = k$ (with $1 \in k$ mapping to the
unit in $W_1$), so similarly, 
$\on{QHH}_{\alpha_0}(W_n^{op}) = W_{n-1}^{op}$. This is the left hand side of
the third map in \eqref{eq:weyl-bimod} in our setting.

\subsubsection{}

Next, we calculate the right hand side of \eqref{eq:weyl-bimod}.

First, we have:
%
\[
\ul{\End}_{\IndCoh^*(\sY^{\leq n})}(\sF_n) \isom 
\ul{\End}_{\IndCoh^*(\sY)}(\lambda_{n,*}^{\IndCoh}(\sF_n))
\]

\noindent by Remark \ref{r:kash-y} (and the assumption that $n>0$).

By \eqref{eq:fund-ses}, we then have:
%
\[
\ul{\End}_{\IndCoh^*(\sY)}\big(\lambda_{n,*}^{\IndCoh}(\sF_n)\big)^{\alpha_0\cdot-} = 
\ul{\Hom}_{\IndCoh^*(\sY)}\big(\lambda_{n,*}^{\IndCoh}(\sF_n),
\lambda_{n-1,*}^{\IndCoh}(\sF_{n-1})\big). 
\]

\noindent Applying \eqref{eq:fund-ses} again, we have:
%
\begin{equation}\label{eq:qh-end}
\source{tates-thesis-de-rham-setting}
\begin{gathered}
\on{QH}_{\alpha_0}\Big(\ul{\End}_{\IndCoh^*(\sY)}\big(\lambda_{n,*}^{\IndCoh}(\sF_n)\big)\Big) 
\coloneqq \\
\ul{\End}_{\IndCoh^*(\sY)}\big(\lambda_{n,*}^{\IndCoh}(\sF_n)\big)^{\alpha_0\cdot-}/-\cdot\alpha_0 =
\ul{\End}_{\IndCoh^*(\sY)}\big(\lambda_{n-1,*}^{\IndCoh}(\sF_{n-1})\big).
\end{gathered}
\end{equation}

By induction, we know that the right hand side identifies with 
$W_{n-1}^{op}$ as an algebra. Below, we will show that the
comparison map from Lemma \ref{l:weyl-bimod} coincides with this
identification.

\subsubsection{}

Let us return to the general setting of \S \ref{sss:qh-1}.

First, observe that $\on{QH}_{\alpha}(A)$ carries a canonical algebra structure.
Indeed, it obviously identifies with:
%
\[
\ul{\End}_{A\mod}(A/-\cdot\alpha) = 
\ul{\End}_{A\mod}(A \underset{k[t]}{\otimes} k).
\]

\noindent (Here $k[t]$ acts on $k$ with $t$ acting by $0$.)
We equip $A^{\alpha\cdot-}/-\cdot \alpha$ with the multiplication 
\emph{opposite} to this one.\footnote{This convention is natural 
because $\ul{\End}_{A\mod}(A) = A^{op}$.}

Suppose now that $A$ and $B$ are (DG) algebras, and that we are given a map:
%
\[
B[t] \coloneqq B \otimes k[t] \to A.
\] 

We can further identify:
%
\[
A \underset{k[t]}{\otimes} k \simeq A \underset{B[t]}{\otimes} B.
\]

\noindent I.e., $A \underset{k[t]}{\otimes} k$ is canonically 
an $(A,B)$-bimodule. Therefore, by construction, we obtain
a canonical morphism: 
%
\[
B \to \on{QH}_{\alpha}(A).
\]

\begin{rem}
\source{tates-thesis-de-rham-setting}
\notready

We note that the diagram:
%
\[
\xymatrix{
B \ar[rr] \ar[d] && \on{QH}_{\alpha}(A)\ar[d] \\
A \ar[rr] && A/-\cdot\alpha
}
\]

\noindent commutes by construction. 

\end{rem}

\subsubsection{}

We apply the above with:
%
\[
A = \ul{\End}_{\IndCoh^*(\sY^{\leq n})}(\sF_n),
\alpha = \alpha_0, B = W_{n-1}^{op}.
\]

We first see that:
%
\[
\on{QH}_{\alpha_0}\Big(\ul{\End}_{\IndCoh^*(\sY)}\big(\lambda_{n,*}^{\IndCoh}(\sF_n)\big)\Big) 
\]

\noindent carries a canonical algebra structure; by construction, 
\eqref{eq:qh-end} is an isomorphism of algebras, where the right hand side
is given the usual algebra structure on endomorphisms. 

Moreover, we deduce that the map:
%
\[
W_{n-1}^{op} \to \on{QH}_{\alpha_0}\Big(\ul{\End}_{\IndCoh^*(\sY)}\big(\lambda_{n,*}^{\IndCoh}(\sF_n)\big)\Big).
\]

\noindent (coming from Lemma \ref{l:weyl-bimod}) is an algebra morphism.

\subsubsection{}

Finally, we claim that the resulting map:
%
\[
W_{n-1}^{op} \to \on{QH}_{\alpha_0}\Big(\ul{\End}_{\IndCoh^*(\sY)}\big(\lambda_{n,*}^{\IndCoh}(\sF_n)\big)\Big) = 
\ul{\End}_{\IndCoh^*(\sY)}\big(\lambda_{n-1,*}^{\IndCoh}(\sF_{n-1})\big)
\]

\noindent coincides with the construction of \S \ref{s:weyl} (for $n-1$ instead of $n$).

Indeed, as this is an algebra map, we can check this on generators of
$W_{n-1}$. There it is essentially tautological from the constructions.

In particular, we see by induction that this map is an isomorphism.
Therefore, the conditions of Lemma \ref{l:weyl-bimod} are satisfied,
so we have proved Proposition \ref{p:ff}.
